% =====================================================================
%  ACADEMIC CV — TEMPLATE FOR US PhD APPLICATIONS
%  Compile on Overleaf with pdfLaTeX. Uses only the same packages as the
%  reference CV, so it will build without extra setup.
%
%  HOW TO USE
%  1. Replace every <<PLACEHOLDER>> with your own information.
%  2. Delete any section you cannot fill honestly. An absent section is
%     invisible; an empty or padded one is not.
%  3. Keep the total to 2 pages. If you overflow, cut projects first,
%     then coursework, then the References section.
%  4. Retarget the Research Interests block for each professor you write
%     to. It is the only part that should change between applications.
% =====================================================================

\documentclass[a4paper,10pt]{article}

\usepackage[left=1.6cm, right=1.6cm, top=1.5cm, bottom=1.5cm]{geometry}
\usepackage{titlesec}
\usepackage{enumitem}
\usepackage{xcolor}
\usepackage{multicol}
\usepackage{fontawesome5}
\usepackage[hidelinks]{hyperref}

% One link colour, used everywhere. Do not mix black and blue links.
\definecolor{linkblue}{HTML}{1A4C8B}
\hypersetup{colorlinks=true, urlcolor=linkblue, linkcolor=linkblue}

\titleformat{\section}{\bfseries\large\scshape}{}{0pt}{}[\titlerule]
\titlespacing{\section}{0pt}{9pt}{5pt}

\setlength{\parindent}{0pt}
\setlist[itemize]{noitemsep, topsep=1pt, leftmargin=1.2em, label=\textbullet}

% Left-aligned entry title with a right-aligned date on the same line.
\newcommand{\entry}[2]{\textbf{#1}\hfill{#2}\\}

\begin{document}
\pagestyle{empty}

% ===================== HEADER =====================
% City, Country format. No street address, no photo, no date of birth,
% no marital status, no national ID. US committees do not want them and
% in some cases are not permitted to consider them.
\begin{center}
    {\fontsize{18}{22}\selectfont\textbf{<<YOUR FULL NAME>>}}\\[4pt]
    <<City>>, <<Country>>\\[3pt]
    \faEnvelope\ \href{mailto:<<permanent.email@gmail.com>>}{<<permanent.email@gmail.com>>} \quad
    \faPhone\ +<<country code>><<phone>> \quad
    \faLinkedin\ \href{https://linkedin.com/in/<<handle>>}{linkedin.com/in/<<handle>>} \quad
    \faGithub\ \href{https://github.com/<<handle>>}{github.com/<<handle>>}
    % Add if you have them:
    % \quad \faGraduationCap\ \href{<<scholar url>>}{Google Scholar}
    % \quad \faGlobe\ \href{<<site url>>}{<<yoursite.github.io>>}
\end{center}
\vspace{4pt}

% ===================== RESEARCH INTERESTS =====================
% 2-3 lines. Name the SUBFIELD and the METHODS, not your personality.
% Rewrite this for every professor: the words here should overlap with
% the words on their lab page. Delete adjectives like "passionate",
% "highly motivated", "deep passion" - they carry no information.
\section{Research Interests}
<<Subfield 1>>; <<Subfield 2>>; <<Method or technique 1>>; <<Method or technique 2>>.
Specific interest in <<the one concrete problem you want to work on, in one clause>>.

% ===================== EDUCATION =====================
\section{Education}
\entry{<<University Name, Spelled Out>>}{<<Mon Year>> -- <<Mon Year / Expected Mon Year>>}
B.Sc. in <<Full Official Degree Name>>\\
<<City>>, <<Country>>\\
CGPA: <<X.XX>> / 4.00 \quad|\quad <<Rank if strong: e.g. Rank 5 of 180>>\\
Thesis: \textit{<<Thesis Title>>} \quad|\quad Advisor: <<Prof. Name>>

% Add earlier degrees only if they are degrees. Do not list high school.

% ===================== RESEARCH EXPERIENCE =====================
% This is the section a PhD advisor reads first. Give it the most space.
% Every bullet: ACTION VERB + WHAT/HOW + RESULT.
% Present tense for ongoing work, past tense for finished work.
% Put a NUMBER in at least half the bullets.
\section{Research Experience}

\entry{<<Position: e.g. Undergraduate Thesis Researcher>>}{<<Mon Year>> -- <<Present>>}
<<Lab or Group Name>>, <<University>> \quad|\quad Advisor: \href{<<faculty page url>>}{<<Prof. Name>>}
\begin{itemize}
  \item <<Verb>> <<what you built or modelled, with the specific method/tool>>, <<result: accuracy, speedup, error, agreement with theory>>
  \item <<Verb>> <<the dataset or simulation campaign>>, <<size: N samples, M sweeps, parameter range>>
  \item <<Verb>> <<the comparison or validation you ran>>, <<what it showed vs. the baseline>>
  \item <<Verb>> <<a concrete artefact: codebase, tool, measurement rig>>, \href{<<repo url>>}{code available}
  \item <<Verb>> <<what you are doing now / what is next>>, targeting submission to <<venue>>
\end{itemize}

\vspace{4pt}
\entry{<<Position: e.g. Research Assistant>>}{<<Mon Year>> -- <<Mon Year>>}
<<Lab or Group Name>>, <<University>> \quad|\quad Supervisor: <<Prof. Name>>
\begin{itemize}
  \item <<Verb>> <<context>>, <<result>>
  \item <<Verb>> <<context>>, <<result>>
\end{itemize}

% ===================== PUBLICATIONS =====================
% Bold your own name. Number the entries so a reader can count them.
% State status honestly. Never list "in preparation" work you have not
% actually started writing.
\section{Publications}
\textbf{Peer-Reviewed Conference Papers}
\begin{enumerate}[leftmargin=1.4em, itemsep=3pt]
  \item \textbf{<<Y. Z. Yourname>>}, <<Co-author>>, and <<Co-author>>,
        ``<<Paper Title>>,'' in \textit{<<Full Conference Name (Acronym)>>},
        <<City, Country>>, <<Year>>, pp.~<<xx--yy>>.
        \href{<<doi url>>}{doi:<<10.xxxx/yyyy>>}
\end{enumerate}

\vspace{4pt}
\textbf{Under Review}
\begin{enumerate}[leftmargin=1.4em, itemsep=3pt, resume]
  \item \textbf{<<Y. Z. Yourname>>} and <<Co-author>>,
        ``<<Paper Title>>,'' submitted to \textit{<<Journal Name>>}, <<Mon Year>>.
\end{enumerate}

% Delete either sub-block if empty. If you have zero publications,
% delete this whole section - an empty Publications heading is worse
% than none, and plenty of admitted PhD students had none.

% ===================== CONFERENCE PRESENTATIONS =====================
% Only if you personally presented. Delete otherwise.
\section{Conference Presentations}
\textbf{<<Yourname, Y. Z.>>} ``<<Talk or Poster Title>>.'' <<Conference Name>>,
<<City, Country>>, <<Day Mon Year>>. <<Oral Presentation / Poster>>.

% ===================== SELECTED PROJECTS =====================
% THREE to FIVE only, chosen to support the research story above.
% A course lab is not a project unless you went beyond the handout.
% If a project has no result worth stating, it does not belong here.
\section{Selected Projects}

\entry{<<Project Name>>}{<<Mon Year>>}
<<One line: what it is and which course/lab, if any>>
\begin{itemize}
  \item <<Verb>> <<what you implemented, with the tool>>, <<measured result>>
  \item <<Verb>> <<what you verified or optimised>>, <<number>>
\end{itemize}
\href{<<repo url>>}{\faGithub\ Repository}

\vspace{5pt}
\entry{<<Project Name>>}{<<Mon Year>>}
<<One line description>>
\begin{itemize}
  \item <<Verb>> <<context>>, <<result>>
  \item <<Verb>> <<context>>, <<result>>
\end{itemize}
\href{<<repo url>>}{\faGithub\ Repository}

\vspace{5pt}
\entry{<<Project Name>>}{<<Mon Year>>}
<<One line description>>
\begin{itemize}
  \item <<Verb>> <<context>>, <<result>>
\end{itemize}
\href{<<repo url>>}{\faGithub\ Repository}

% ===================== TECHNICAL SKILLS =====================
% Subcategorise. Add proficiency levels only if they genuinely vary.
% Do not list soft skills. Do not list Microsoft Word.
% Only list a tool you could be questioned about in an interview.
\section{Technical Skills}
\textbf{Programming:} <<Advanced -- Python, MATLAB; Intermediate -- C/C++, SystemVerilog>>\\
\textbf{Simulation \& Design:} <<Tool, Tool, Tool>>\\
\textbf{Experimental / Lab:} <<Instrument, technique, fabrication step>>\\
\textbf{Frameworks \& Libraries:} <<PyTorch, NumPy, SciPy, ...>>\\
\textbf{Languages:} English (<<IELTS X.X / TOEFL XXX>>), <<Native language>> (Native)

% ===================== TEST SCORES =====================
% US committees look for these. Put real numbers or delete the section.
% Do not write "to be taken" - that reads as not taken.
\section{Test Scores}
\textbf{IELTS / TOEFL iBT:} <<Overall X.X>> \quad (L <<x>>, R <<x>>, W <<x>>, S <<x>>) \hfill <<Mon Year>>\\
\textbf{GRE General:} Quantitative <<xxx>>, Verbal <<xxx>>, Analytical Writing <<x.x>> \hfill <<Mon Year>>

% ===================== TEACHING EXPERIENCE =====================
% Matters because most US PhD funding runs through TA-ships. Include
% tutoring and grading if that is what you have.
\section{Teaching Experience}
\entry{<<Teaching Assistant, Course Name>>}{<<Mon Year>> -- <<Mon Year>>}
<<Department>>, <<University>>
\begin{itemize}
  \item <<Verb>> <<lab sections / recitations>> for <<N>> students, covering <<topics>>
  \item <<Verb>> <<graded, designed problem sets, held office hours>>, <<outcome>>
\end{itemize}

% ===================== HONORS AND AWARDS =====================
% Reverse chronological. Add one clause of context for anything a US
% reader cannot calibrate ("top 1% of 20,000 candidates nationally").
% Nothing from high school except a truly national-level distinction.
\section{Honors and Awards}
\begin{itemize}
  \item \textbf{<<Award Name>>}, <<Awarding Body>> \hfill <<Mon Year>>\\
        <<One clause of scale/selectivity, e.g. 1 of 12 teams selected from 140 entries worldwide>>
  \item \textbf{<<Award Name>>}, <<Awarding Body>> \hfill <<Year>>
  \item \textbf{<<Scholarship Name>>}, <<Awarding Body>> -- <<selection basis>> \hfill <<Year>>
\end{itemize}

% ===================== LEADERSHIP AND SERVICE =====================
% Keep short. Departmental, professional-society, or peer-review service.
\section{Leadership and Service}
\begin{itemize}
  \item <<Role>>, <<Organisation>> \hfill <<Mon Year>> -- <<Mon Year>>
  \item <<Role>>, <<Organisation>> \hfill <<Mon Year>> -- <<Mon Year>>
\end{itemize}

% ===================== RELEVANT COURSEWORK =====================
% SIX at most, all pointed at the lab you are applying to. Include the
% grade only if it is uniformly strong. Put this last: it is the least
% informative section on the page. Delete it if you are over 2 pages.
\section{Relevant Coursework}
\begin{multicols}{2}
\begin{itemize}
  \item <<Course>> (<<A / 4.00>>)
  \item <<Course>> (<<A / 4.00>>)
  \item <<Course>> (<<A / 4.00>>)
  \item <<Course>> (<<A / 4.00>>)
  \item <<Course>> (<<A / 4.00>>)
  \item <<Course>> (<<A / 4.00>>)
\end{itemize}
\end{multicols}

% ===================== REFERENCES =====================
% OPTIONAL, and the first thing to cut for space. Your recommenders are
% submitted through the application portal, so this is usually redundant
% for a formal application. It IS worth keeping in the version you email
% directly to a professor, because it shows who can vouch for you.
% If you keep it: ask permission first, and use their institutional email.
\section{References}
\begin{minipage}[t]{0.48\textwidth}
\href{<<faculty page url>>}{<<Prof. Full Name>>}\\
<<Title>>, <<Department>>\\
<<University>>\\
\href{mailto:<<email>>}{<<email>>}
\end{minipage}
\begin{minipage}[t]{0.48\textwidth}
\href{<<faculty page url>>}{<<Prof. Full Name>>}\\
<<Title>>, <<Department>>\\
<<University>>\\
\href{mailto:<<email>>}{<<email>>}
\end{minipage}

\end{document}
